\subsection{Physical Properties Controlling Thermoelastic Wave Propagation}

Thermoelastic wave propagation in rocks is controlled by a group of fundamental physical properties that determine how a geological medium stores mass, resists deformation, conducts heat, and develops thermally induced strain. In real rocks, these properties do not act independently. Density controls inertia, elastic stiffness controls resistance to strain, porosity modifies both the mechanical framework and thermal transport pathways, and thermal parameters govern the rate at which temperature disturbances are absorbed and redistributed. Because rocks are natural materials rather than ideal solids, each of these properties depends on lithology, mineral composition, pore structure, fracture density, saturation, pressure, temperature, and the scale of observation. For this reason, thermoelastic behavior in geological media is best understood as the coupled response of a heterogeneous solid rather than as the consequence of a single material constant \cite{mavko2009,jaeger2007,robertson1988,turcotte2014}. % Verified with web sources: 

Density, denoted by $\rho$, is the mass per unit volume of a material and is commonly expressed in $\mathrm{kg,m^{-3}}$. In rock physics, it is often useful to distinguish between grain or matrix density and bulk density, because the latter includes the effect of pore space and the material occupying the pores. Density is important for thermoelastic wave propagation because it represents inertia: the greater the density of the medium, the greater the resistance of that medium to acceleration under an applied stress. In an isotropic elastic approximation, this role appears directly in the familiar expressions \[ V_P=\sqrt{\frac{K+\frac{4}{3}G}{\rho}}, \qquad V_S=\sqrt{\frac{G}{\rho}}, \] where $V_P$ and $V_S$ are the compressional and shear wave velocities, and $K$ and $G$ are the bulk and shear moduli, respectively. These relations show that wave speed does not depend on stiffness alone; it depends on stiffness relative to density. A rock may therefore be dense but still transmit waves rapidly if its elastic moduli are sufficiently high, whereas porous or weakly consolidated rocks commonly exhibit lower velocities because reduced stiffness is not fully compensated by lower density \cite{aki2002,mavko2009}. % Verified with web sources: 

Porosity is the fraction of the bulk rock volume occupied by pores, voids, or cracks. A useful distinction is made between total porosity and effective porosity. Total porosity refers to all void space within the rock, whereas effective porosity refers to the interconnected pore space that can participate in fluid transmission. This distinction is particularly important in geological materials since some of the pores may be isolated, while others may be interconnected through grain boundaries, microcracks, vugs or fracture systems. The physical effect of porosity depends not only on the pore volume, but also on pore shape, aspect ratio, size distribution, and connectivity. Thin compliant cracks may strongly reduce effective stiffness even when their volume fraction is small, while rounded isolated pores may have a different mechanical influence. Saturation introduces an additional level of complexity: the presence of water, gas, or mixed fluids changes bulk density, fluid compressibility, thermal contact within the pore space, and the way elastic disturbances interact with the solid framework. As a result, porosity influences wave propagation through several coupled mechanisms, including reduced stiffness, modified inertia, altered heat transfer, and changes in the distribution of thermal stress \cite{wolff1982,mavko2009,jaeger2007}. % Verified with web sources: 

Elastic stiffness expresses the resistance of a rock to deformation under applied stress. In a general sense, a stiffer rock requires a larger stress to produce a given strain, whereas a more compliant rock deforms more readily. In isotropic linear elasticity, stiffness is commonly characterized by Young's modulus $E$, the bulk modulus $K$, and the shear modulus $G$. Young's modulus describes resistance to axial deformation, the bulk modulus describes resistance to volumetric compression, and the shear modulus describes resistance to changes in shape. These moduli are directly relevant to wave propagation because elastic waves are disturbances that move through a medium by repeated local storage and release of elastic energy. Consequently, rocks with higher effective stiffness generally transmit elastic disturbances faster than rocks with lower stiffness, provided density is comparable. In real geological media, however, the operative stiffness is an effective property of the rock mass and is influenced by microcracks, cementation, grain contacts, anisotropy, pressure, and temperature, so the observed wave response reflects the structure of the medium rather than the mineral grains alone \cite{mavko2009,jaeger2007,aki2002}. % Verified with web sources: 

Poisson's ratio, usually denoted by $\nu$, describes the relation between longitudinal and lateral strain during elastic deformation. Conceptually, it measures how much a material tends to contract laterally when compressed axially, or expand laterally when stretched axially. In rock mechanics and seismology, this parameter is useful because it reflects the balance between volumetric and shear deformation. Higher Poisson’s ratio of the material corresponds to more lateral deformation under uniaxial loading and lower value indicates relatively less lateral strain. For wave propagation, Poisson's ratio is relevant because it is linked to the relative behavior of compressional and shear disturbances, and therefore to the ratio between $V_P$ and $V_S$. In rocks, Poisson's ratio is not merely a fixed intrinsic constant. It may vary with porosity, crack geometry, mineral composition, saturation, and anisotropy, so changes in this parameter may signal changes in the internal structure or stress state of the medium \cite{mavko2009,jaeger2007,aki2002}. % Verified with web sources: 

The thermal expansion coefficient, denoted here by $\alpha$, characterizes the strain produced by a unit change in temperature and is commonly expressed in $\mathrm{K^{-1}}$. In the simplest isotropic form, thermal strain may be written approximately as $\varepsilon_{th}=\alpha \Delta T$, which means that heating tends to produce expansion and cooling tends to produce contraction. For rocks, this property is particularly important because geological materials are commonly polymineralic. Different minerals may expand by different amounts, and in anisotropic crystals the amount of expansion may also depend on crystallographic direction. As a result, temperature change may generate internal incompatibilities of strain between neighboring grains even when the rock is free of external mechanical loading. Where deformation is constrained by surrounding material, bedding, grain contacts, or fractures, these incompatibilities may produce local thermal stress and microcracking. In this way, thermal expansion is not only a thermal property but also a key link between temperature change and mechanical response in rocks \cite{robertson1988,jaeger2007}. % Verified with web sources: 

Thermal conductivity, denoted here by $k$, measures the ability of a material to transmit heat by conduction and is commonly expressed in $\mathrm{W,m^{-1},K^{-1}}$. In geological media, conductive heat transfer is the principal mechanism by which temperature disturbances spread through the solid framework. Thermal conductivity depends strongly on mineral composition because different minerals conduct heat with different efficiency. It is also influenced by pores and cracks, which interrupt heat paths through the solid; by the shape and connectivity of pores; and by the nature of the pore-filling fluid. Air-filled pores generally reduce conductivity more strongly than water-filled pores, because air is a poorer conductor than water and most rock-forming minerals. In addition, thermal conductivity may be anisotropic in layered, foliated, or strongly fractured rocks, so heat may move more readily in one direction than in another. This directional aspect is important for thermoelastic problems because the spatial evolution of temperature governs the spatial evolution of thermal strain and stress \cite{robertson1988,turcotte2014}. % Verified with web sources: 

Specific heat capacity at constant pressure, written as $C_p$, is the amount of heat required to raise the temperature of unit mass of material by one degree and is commonly expressed in $\mathrm{J,kg^{-1},K^{-1}}$. In physical terms, it measures the heat-storage capacity of the material. The rock with higher $C_p$ can absorb more thermal energy before its temperature changes by a certain amount, while the rock with lower $C_p$ reacts more quickly to the same heat input. This property is important for thermoelastic wave propagation as it controls the rate of evolution of temperature during transient heating or cooling. Multiplying by density gives the volumetric heat capacity $\rho C_p$ which is often more directly relevant in geological problems as it expresses heat storage per unit volume of rock. The temperature field generated by a thermal disturbance therefore depends not only on how efficiently heat is conducted, but also on how much energy must be stored to change the temperature of the medium \cite{robertson1988,turcotte2014}. % Verified with web sources: 

A closely related property is thermal diffusivity, denoted here by $\alpha_T$. Using the notation adopted in this thesis, it is defined as \[ \alpha_T=\frac{k}{\rho C_p}, \] where $k$ is thermal conductivity, $\rho$ is density, and $C_p$ is specific heat capacity. Thermal diffusivity describes how rapidly a temperature disturbance spreads through the rock. A high $\alpha_T$ means that heat is redistributed relatively quickly, while a low $\alpha_T$ means that thermal gradients persist for a longer time. This distinction is important in thermoelasticity because sharp or persistent thermal gradients can produce stronger local thermal stress than rapidly equilibrating temperature fields. It is also necessary to emphasize that thermal diffusivity $\alpha_T$ must not be confused with the thermal expansion coefficient $\alpha$. Although the symbols are similar, they describe entirely different physical effects: $\alpha$ governs thermal strain, whereas $\alpha_T$ governs the rate of temperature diffusion \cite{robertson1988,turcotte2014}. % Verified with web sources: 

In real geological media, the properties discussed above are not strictly constant. Pressure commonly affects rocks by closing compliant cracks and reducing pore volume, which tends to increase effective stiffness and elastic wave velocity. Temperature may reduce elastic moduli, modify thermal transport, and promote thermally induced microdamage. Saturation changes both mechanical and thermal behavior because pore fluids alter density, effective compressibility, and conductive heat transfer. Finally, scale is crucial: values measured on small laboratory specimens may differ from the effective properties of a field-scale rock mass that contains bedding planes, joints, faults, weathered zones, and fracture networks. For this reason, any theoretical treatment of thermoelastic wave propagation must recognize that physical properties in rocks are conditional properties, valid only for a given range of temperature, pressure, saturation state, fabric, and observational scale \cite{mavko2009,jaeger2007,robertson1988,wolff1982,turcotte2014}. % Verified with web sources: 

Taken together, density, porosity, elastic stiffness, Poisson's ratio, thermal expansion, thermal conductivity, specific heat capacity, and thermal diffusivity define the physical framework within which thermoelastic disturbances evolve in rock. They determine how fast a wave travels, how strongly the medium resists deformation, how temperature changes are stored and transmitted, and how thermal strain is converted into stress. The following sections therefore build on these properties in order to examine elastic deformation, stress--strain relations, heat conduction, and the coupled thermoelastic behavior of geological media \cite{mavko2009,aki2002,jaeger2007,robertson1988}. % Verified with web sources: 